% SPDX-FileCopyrightText: Copyright © 2026 Idiap Research Institute <contact@idiap.ch>
%
% SPDX-FileContributor: Andre Freitas <andre.freitas@idiap.ch>
% Neuro-symbolic AI Group
%
% SPDX-FileContributor: Danilo Gusicuma <danilo.gusicuma@idiap.ch>
% Neuro-symbolic AI Group
%
% SPDX-License-Identifier: MIT
\documentclass[11pt]{article}

\usepackage[a4paper,margin=1in]{geometry}
\usepackage{amsmath,amssymb,amsthm,mathtools}
\usepackage{enumitem}
\usepackage{hyperref}
\usepackage{booktabs}
\usepackage{array}
\usepackage{xcolor}

\hypersetup{
    colorlinks=true,
    linkcolor=blue!60!black,
    citecolor=blue!60!black,
    urlcolor=blue!60!black
}

\newtheorem{definition}{Definition}
\newtheorem{proposition}{Proposition}
\newtheorem{remark}{Remark}
\newtheorem{example}{Example}

\title{A Hegelian Logic of Mathematical Concepts:\\
Determinate Negation, Closure Failure, and Conceptual Enrichment}
\author{}
\date{}

\begin{document}

\maketitle

\begin{abstract}
This note develops a precise interpretation of Hegelian logic for mathematical concepts. The central claim is not that mathematical proof is dialectical in the ordinary sense, nor that formal inference should be replaced by speculative philosophy. Rather, Hegelian logic can be used as a theory of \emph{conceptual development}: mathematical concepts often become more expressive when their initial forms encounter determinate limits, such as failures of closure, solvability, completeness, invariance, expressivity, or structural adequacy. The new concept cancels the inadequacy of the previous one, preserves it as a special case, and elevates it into a richer domain. This is formalized as a mathematical analogue of \emph{Aufhebung}: preservation-through-transformation.
\end{abstract}

\section{Introduction}

Hegelian logic is usually understood as a logic of conceptual movement rather than a logic of static classification. A concept is not fully grasped by assigning it a fixed definition; it must also be understood through the tensions, limits, and transformations that arise when the concept is developed rigorously. In mathematical contexts, this perspective can be made unusually concrete.

Mathematical concepts frequently begin in restricted forms. Natural numbers express counting; functions are initially understood as formulas; equality appears to mean strict identity; infinity appears as endless continuation. Yet each of these concepts encounters precise limitations. Subtraction is not closed in the natural numbers. Formula-based functions are too narrow. Strict equality is too rigid for structural mathematics. Potential infinity is insufficient for modern set theory.

The Hegelian reading proposed here is that these are not merely external historical accidents. They are conceptually productive failures. A concept becomes richer when its limits become explicit. The resulting development has the form:
\[
C_0 \longrightarrow L_0 \longrightarrow C_1,
\]
where $C_0$ is an initial concept, $L_0$ is a determinate limitation, and $C_1$ is a richer concept that preserves and transforms $C_0$.

This is a mathematical analogue of Hegel's notion of \emph{Aufhebung}: a movement that simultaneously cancels, preserves, and elevates. The earlier concept is not simply refuted. It is retained as a special case within a more adequate structure.

\section{The General Schema}

Let $C$ denote a mathematical concept or conceptual domain. A Hegelian analysis does not ask only what $C$ is, but also what $C$ cannot yet do. The developmental pressure on $C$ arises when a mathematically meaningful operation, relation, invariant, or theorem-schema cannot be expressed or satisfied within $C$.

\begin{definition}[Conceptual limitation]
A \emph{conceptual limitation} of a mathematical concept $C$ is a structured failure
\[
L(C,O,P)
\]
where $O$ is an operation, construction, or interpretive demand, and $P$ is a property expected of $C$ but not satisfied within its current formulation.
\end{definition}

Typical properties include closure, solvability, completeness, invariance, expressive adequacy, and structural naturality.

\begin{definition}[Determinate negation]
A limitation $L(C,O,P)$ is a \emph{determinate negation} of $C$ when it does not merely reject $C$, but identifies a specific missing structure whose addition generates a more expressive concept $C'$.
\end{definition}

Thus, determinate negation differs from abstract negation. Abstract negation says only that a concept fails. Determinate negation says how it fails, and therefore in which direction the concept must be transformed.

\begin{definition}[Mathematical Aufhebung]
Given a concept $C$ and a determinate limitation $L$, an \emph{Aufhebung} of $C$ with respect to $L$ is a concept $C'$ together with an embedding or interpretation
\[
\iota : C \hookrightarrow C'
\]
such that:
\begin{enumerate}[label=(\roman*)]
    \item $C'$ resolves, generalizes, or relocates the limitation $L$;
    \item $C$ is preserved inside $C'$ as a substructure, special case, quotient, approximation, or local model;
    \item $C'$ makes explicit a structure that was only implicit or impossible in $C$.
\end{enumerate}
\end{definition}

This allows the Hegelian notion of cancellation-preservation-elevation to be given a more mathematical form:
\[
C' = \operatorname{Aufhebung}(C,L).
\]

\section{Number Systems as Dialectical Development}

The classical expansion of number systems is the clearest example.

\subsection{From Natural Numbers to Integers}

The natural numbers
\[
\mathbb{N} = \{0,1,2,3,\dots\}
\]
express counting. They are closed under addition and multiplication, but not under subtraction. For instance,
\[
3-5 \notin \mathbb{N}.
\]
This is not an arbitrary defect. It is a determinate failure: the operation of subtraction demands additive inverses. The concept of number is therefore enlarged to the integers:
\[
\mathbb{Z} = \{\dots,-2,-1,0,1,2,\dots\}.
\]
The natural numbers are preserved inside the integers:
\[
\mathbb{N} \subset \mathbb{Z}.
\]
Thus $\mathbb{Z}$ cancels the insufficiency of $\mathbb{N}$ with respect to subtraction, while preserving $\mathbb{N}$ as the nonnegative part of a richer additive structure.

\subsection{From Integers to Rationals}

The integers are not closed under division. For example,
\[
3/5 \notin \mathbb{Z}.
\]
The determinate failure concerns multiplicative inverses. The concept is enlarged to the rational numbers:
\[
\mathbb{Q} = \left\{\frac{a}{b} : a,b \in \mathbb{Z},\ b\neq 0\right\}.
\]
Again, the previous concept is not abandoned:
\[
\mathbb{Z} \subset \mathbb{Q}.
\]

\subsection{From Rationals to Reals}

The rationals are algebraically rich, but incomplete as an ordered metric field. There are Cauchy sequences of rational numbers that do not converge to rational numbers. Geometrically, the diagonal of the unit square has length
\[
\sqrt{2},
\]
which is not rational. Thus the rationals fail to capture continuity and metric completeness.

The reals $\mathbb{R}$ complete $\mathbb{Q}$:
\[
\mathbb{Q} \subset \mathbb{R}.
\]
The earlier structure is preserved, but its limitation is overcome by embedding it into a complete ordered field.

\subsection{From Reals to Complex Numbers}

The real numbers are complete as an ordered field, but they are not algebraically closed. The equation
\[
x^2+1=0
\]
has no solution in $\mathbb{R}$. The complex numbers introduce an element $i$ such that
\[
i^2=-1.
\]
Thus:
\[
\mathbb{R} \subset \mathbb{C}.
\]
The real numbers are preserved as the real axis within the complex plane. The concept of number is elevated from ordered magnitude to algebraically closed field-like structure.

\subsection{Summary}

The developmental chain is:
\[
\mathbb{N}
\longrightarrow
\mathbb{Z}
\longrightarrow
\mathbb{Q}
\longrightarrow
\mathbb{R}
\longrightarrow
\mathbb{C}.
\]
Each transition is generated by a determinate failure:

\begin{center}
\begin{tabular}{lll}
\toprule
Initial concept & Limitation & Enriched concept \\
\midrule
$\mathbb{N}$ & not closed under subtraction & $\mathbb{Z}$ \\
$\mathbb{Z}$ & not closed under division & $\mathbb{Q}$ \\
$\mathbb{Q}$ & not complete & $\mathbb{R}$ \\
$\mathbb{R}$ & not algebraically closed & $\mathbb{C}$ \\
\bottomrule
\end{tabular}
\end{center}

This is a paradigmatic mathematical instance of Aufhebung: each new domain negates the limitation of the previous one while preserving the previous one as a subdomain.

\section{Geometry: From Euclidean Space to Manifolds}

Euclidean geometry begins with an apparently self-evident spatial structure: points, lines, planes, distances, angles, and the parallel postulate. Yet the parallel postulate introduced a conceptual tension: was Euclidean space the only possible geometry, or merely one geometry among others?

The development of non-Euclidean geometries shows that Euclidean geometry was not false, but partial. Hyperbolic and elliptic geometries do not merely reject Euclidean geometry. They place it inside a broader space of possible geometrical structures.

The trajectory can be schematized as:
\[
\text{Euclidean geometry}
\longrightarrow
\text{parallel-postulate problem}
\longrightarrow
\text{non-Euclidean geometries}.
\]

Later, differential geometry further generalizes this movement:
\[
\text{Euclidean spaces}
\longrightarrow
\text{topological spaces}
\longrightarrow
\text{smooth manifolds}
\longrightarrow
\text{Riemannian manifolds}.
\]

Here the earlier concept is preserved locally. A smooth manifold is locally Euclidean, even though its global structure may be curved or topologically nontrivial. In Hegelian terms, Euclidean geometry is not destroyed; it is aufgehoben as a local model within a more general geometry.

\section{Functions: From Formula to Morphism}

The concept of function initially appears as a formula:
\[
f(x)=x^2+3x+1.
\]
This is a powerful but restricted conception. Many mathematically important functions are not naturally given by elementary formulas. The concept is therefore expanded to that of a mapping:
\[
f:X\to Y.
\]
A function becomes any rule assigning to each input in $X$ exactly one output in $Y$.

Modern mathematics then further transforms the concept. In algebra, topology, and category theory, what matters is often not arbitrary input-output assignment, but structure preservation. A group homomorphism preserves group operations; a continuous map preserves topological structure; a smooth map preserves differentiable structure.

The conceptual development is:
\[
\text{formula}
\longrightarrow
\text{mapping}
\longrightarrow
\text{structure-preserving map}
\longrightarrow
\text{morphism}
\longrightarrow
\text{functor}.
\]

A functor preserves not just structure within objects, but the structure of an entire category:
\[
F:\mathcal{C}\to\mathcal{D}.
\]
It maps objects to objects and morphisms to morphisms while preserving identities and composition. Thus, the original formula-concept of function is preserved as a special case within a much richer theory of structure-preserving transformation.

\section{Equality, Isomorphism, and Structural Sameness}

At first, equality appears to be strict identity:
\[
a=b.
\]
But mathematical practice shows that strict identity is often too rigid. Two groups may not be identical as sets, but may have the same algebraic structure. Two topological spaces may be different presentations of the same topological form. Two categories may be equivalent without being literally equal.

Thus mathematics develops a family of sameness relations:
\[
=,\quad \cong,\quad \simeq,\quad \equiv,\quad \sim,\quad \approx.
\]
These include equality, isomorphism, equivalence, homotopy equivalence, congruence, and approximation.

The dialectical structure is:
\[
\text{strict identity}
\longrightarrow
\text{difference of representation}
\longrightarrow
\text{structural equivalence}.
\]

This is one of the most important conceptual developments in modern mathematics. The relevant object is not always the object as given, but the object considered up to the appropriate equivalence relation. In this sense, equality is aufgehoben into invariance under transformation.

\section{Infinity: From Endlessness to Structured Hierarchy}

Infinity first appears as potential infinity: an endless process, such as the sequence
\[
1,2,3,4,\dots.
\]
One can always continue. Yet modern mathematics treats infinite collections as actual objects, such as
\[
\mathbb{N}, \quad \mathbb{R}, \quad \mathcal{P}(\mathbb{N}).
\]
Set theory then reveals that infinity is not homogeneous. There are different cardinalities of infinite sets:
\[
|\mathbb{N}| < |\mathbb{R}|.
\]
Thus infinity develops from a negative notion, ``without end,'' into a structured hierarchy of sizes and principles.

The trajectory is:
\[
\text{potential infinity}
\longrightarrow
\text{actual infinite set}
\longrightarrow
\text{cardinal hierarchy}
\longrightarrow
\text{large-cardinal principles}.
\]

The initial idea of endlessness is preserved, but it is transformed into a theory of infinite structure.

\section{Types of Mathematical Negation}

The Hegelian notion of negation can be made precise by identifying mathematical failure modes. These are not vague contradictions, but determinate constraints that force conceptual enrichment.

\begin{center}
\begin{tabular}{p{0.25\linewidth}p{0.32\linewidth}p{0.30\linewidth}}
\toprule
Type of limitation & Example & Enriched concept \\
\midrule
Closure failure & $3-5\notin\mathbb{N}$ & Integers \\
Solvability failure & $x^2+1=0$ has no real solution & Complex numbers \\
Completeness failure & Cauchy sequences in $\mathbb{Q}$ need not converge in $\mathbb{Q}$ & Real numbers \\
Invariance failure & Coordinates vary while geometry remains & Differential geometry \\
Expressivity failure & Formula view of function is too narrow & Set-theoretic functions \\
Rigidity failure & Equality is too strict for structural sameness & Isomorphism and equivalence \\
Foundational failure & Naive comprehension yields paradoxes & Axiomatic set theory or type theory \\
Representation failure & Objects depend on presentation & Category-theoretic abstraction \\
\bottomrule
\end{tabular}
\end{center}

These failure modes make Hegelian contradiction mathematically intelligible. A contradiction is not necessarily a literal violation of classical logic. It is often a structural instability in the concept's initial form.

\section{Formal Inference and Conceptual Development}

It is important to distinguish two levels:
\[
\text{formal inference} \neq \text{conceptual development}.
\]

Formal inference operates inside a fixed system. Given axioms, definitions, and rules, one proves theorems. Conceptual development changes the space of possible axioms, definitions, objects, and invariants.

Hegelian logic is most useful at the second level. It helps explain how mathematical objects and frameworks emerge when existing concepts become insufficient. Proofs are not replaced by dialectics; rather, dialectics describes the transformation of the conceptual conditions under which proofs become possible.

\begin{proposition}[Conceptual enrichment pattern]
Many mathematical concept developments can be represented as triples
\[
(C,L,C'),
\]
where $C$ is an initial concept, $L$ is a determinate limitation of $C$, and $C'$ is an enriched concept satisfying a preservation condition
\[
\iota:C\hookrightarrow C'
\]
while resolving or relocating $L$.
\end{proposition}

\begin{remark}
The preservation condition need not always be literal subset inclusion. It may be an embedding, faithful functor, quotient relation, equivalence class construction, local approximation, universal property, or conservative extension.
\end{remark}

\section{Toward a Theory of Mathematical Aufhebung}

A more systematic theory could classify concept-transformations according to the kind of preservation map involved.

\begin{center}
\begin{tabular}{ll}
\toprule
Preservation mode & Mathematical form \\
\midrule
Subset preservation & $\mathbb{N}\subset\mathbb{Z}\subset\mathbb{Q}\subset\mathbb{R}\subset\mathbb{C}$ \\
Embedding & A structure is represented inside a richer structure \\
Completion & $\mathbb{Q}$ becomes dense in $\mathbb{R}$ \\
Quotienting & Objects are identified under an equivalence relation \\
Localization & A structure becomes valid locally rather than globally \\
Conservative extension & New language is added without changing old theorems \\
Categorification & Set-level concepts are lifted to object-and-morphism structures \\
\bottomrule
\end{tabular}
\end{center}

This suggests the following abstract schema:
\[
\operatorname{Aufhebung}(C,L)
= (C',\iota,\rho),
\]
where $\iota:C\to C'$ is a preservation map and $\rho$ is a resolution mechanism for the limitation $L$.

Different mathematical developments instantiate different forms of $\rho$:
\begin{itemize}
    \item adding inverses;
    \item adding limits;
    \item adding solutions;
    \item weakening equality to equivalence;
    \item replacing global structure with local structure;
    \item replacing objects with transformations;
    \item replacing presentations with invariants.
\end{itemize}

\section{Conclusion}

A Hegelian logic of mathematical concepts treats mathematics as a history of conceptual enrichment. Mathematical concepts do not merely accumulate definitions. They transform when their limits become determinate. Natural numbers become integers because subtraction exposes their incompleteness. Rationals become reals because metric completeness exposes their insufficiency. Reals become complex numbers because algebraic solvability exceeds the real domain. Equality becomes isomorphism because strict identity is too rigid for structural mathematics. Functions become morphisms because structure-preservation becomes more central than formulaic expression.

The strongest interpretation is therefore not that proofs are dialectical, but that the evolution of mathematical concepts often follows a dialectical pattern. Initial immediacy encounters determinate negation; the resulting concept preserves the earlier one while transforming its meaning. In mathematical terms:
\[
C_{n+1}=\operatorname{Aufhebung}(C_n,L_n),
\]
where $L_n$ is the precise limitation that forces enrichment. This provides a rigorous bridge between Hegelian logic and the development of mathematical concepts.

\end{document}
